\documentclass[a4paper]{arrowhead}

\usepackage[yyyymmdd]{datetime}
\usepackage{etoolbox}
\usepackage[utf8]{inputenc}
\usepackage[T1]{fontenc}
\usepackage{multirow}
\usepackage{longtable}
\usepackage{tikz}
\usepackage{adjustbox}
\AtBeginEnvironment{longtable}{\small}
\BeforeBeginEnvironment{tikzpicture}{\begin{adjustbox}{max width=\textwidth}}
\AfterEndEnvironment{tikzpicture}{\end{adjustbox}}
\usetikzlibrary{positioning,shapes.multipart,arrows.meta}

\renewcommand{\dateseparator}{-}
\setlength{\emergencystretch}{3em}
\lstset{language=HTTP,breakatwhitespace=false,columns=fullflexible}

%% Special references
\newcommand{\mref}[2][]{{\textcolor{ArrowheadPurple}{\hyperref[sec:model:#2]{\ifstrempty{#1}{#2}{#1}}}}}
\newcommand{\oref}[1]{{\textcolor{ArrowheadBlue}{\hyperref[sec:operations:#1]{#1}}}}
\newcommand{\pdef}[1]{{\textcolor{ArrowheadGrey}{#1\label{sec:model:primitives:#1}\label{sec:model:primitives:#1s}\label{sec:model:primitives:#1es}}}}
\newcommand{\pref}[1]{{\textcolor{ArrowheadGrey}{\hyperref[sec:model:primitives:#1]{#1}}}}

% as template_sd.tex; the input and output arguments are passed already formatted (\mref or \pref)
\newrobustcmd\osubsection[3]{
  \addtocounter{subsection}{1}
  \addcontentsline{toc}{subsection}{\protect\numberline{\thesubsection}operation \textcolor{ArrowheadBlue}{#1}}
  \subsection*{
    \thesubsection\quad
    operation
    \textcolor{ArrowheadBlue}{#1}
    (\notblank{#2}{#2}{})
    \notblank{#3}{: #3}{}
  }
  \label{sec:operations:#1}
}
\newrobustcmd\msubsection[2]{
  \addtocounter{subsection}{1}
  \addcontentsline{toc}{subsection}{\protect\numberline{\thesubsection}#1 \textcolor{ArrowheadPurple}{#2}}
  \subsection*{\thesubsection\quad#1 \textcolor{ArrowheadPurple}{#2}}
  \label{sec:model:#2} \label{sec:model:#2es}
}
%%

\begin{document}
\sloppy

%% Arrowhead Document Properties
\ArrowheadTitle{systemDiscovery}
\ArrowheadServiceID{systemDiscovery}
\ArrowheadType{Service Description}
\ArrowheadTypeShort{SD}
\ArrowheadVersion{5.2.2}
\ArrowheadDate{2026-10-22}
\ArrowheadAuthor{Jerker Delsing}
\ArrowheadStatus{RELEASE CANDIDATE}
\ArrowheadContact{jerker.delsing@ltu.se}
\ArrowheadFooter{\href{www.arrowhead.eu}{www.arrowhead.eu}}
\ArrowheadSetup
%%

%% Front Page
\begin{center}
  \vspace*{1cm}
  \huge{\arrowtitle}

  \vspace*{0.2cm}
  \LARGE{\arrowtype}
  \vspace*{1cm}

  \Large{Service ID: \textit{"\arrowid"}}

  \vspace*{\fill}

  % Front Page Image

  \vspace*{1cm}
  \vspace*{\fill}

  % Front Page Abstract
  \begin{abstract}
    This document describes the \texttt{systemDiscovery} service, version 1.0.0, produced by the ServiceRegistry core system of Arrowhead 5.2.2.
  \end{abstract}

  \vspace*{1cm}
\end{center}

\newpage
%%

%% Table of Contents
\tableofcontents
\newpage
%%

\section{Overview}
\label{sec:overview}

This document describes the \texttt{systemDiscovery} service, which provides registration and lookup of the application, support and core systems of the local cloud, with their addresses, version, metadata and hosting device.
It is produced by the ServiceRegistry system \cite{sr-sysd} in service version 1.0.0.

The rest of this document is organized as follows.
In Section \ref{sec:operations}, we describe the abstract message operations provided by the service.
In Section \ref{sec:model}, we end the document by presenting the data types used by the mentioned operations.
The concrete interface is specified in the HTTP/TLS/JSON IDD \cite{sr-systemDiscovery-http-idd}. The MQTT binding is specified in the MQTT/TLS/JSON IDD \cite{sr-systemDiscovery-mqtt-idd}.

\subsection{Significant Prior Art}

Corresponds to the system part of the Arrowhead 4.x ServiceRegistry \texttt{register}/\texttt{unregister} operations and the 4.4 SystemRegistry.

\subsection{How This Service Is Meant to Be Used}

A system registers itself with \texttt{register} before it registers services; the system name is taken from its authenticated identity, not from the request. \texttt{revoke} removes the requester and, with it, its service instances. \texttt{lookup} finds systems by name, address, version, device or metadata; it is used by Consumer\allowbreak{}Authorization to evaluate metadata-based policies.

\subsection{Important Delimitations and Dependencies}
\label{sec:delimitations}

A system can only register and revoke itself. Registering another system requires \texttt{service\allowbreak{}Registry\allowbreak{}Management}. The service does not authenticate systems; the identity comes from the authentication policy of the ServiceRegistry (certificate, identity token or declared name).

\newpage

\section{Service Interface}
\label{sec:operations}

This section describes the interfaces to the \texttt{systemDiscovery} service.
In particular, each subsection names an abstract operation, an input type and an output type, in that order.
The input type is named inside parentheses, while the output type is preceded by a colon.
Input and output types are only denoted when accepted or returned, respectively, by the interface in question.

All abstract data types named in this section are defined in Section \ref{sec:model}.

\begin{figure}[ht!]
  \centering
  \begin{tikzpicture}[font=\small]
    \node[draw=ArrowheadBlue, very thick, rectangle split, rectangle split parts=2, align=center,
          rectangle split part fill={ArrowheadBlue!12,white}, text width=9cm, inner sep=6pt] (b)
      {\guillemotleft{}Service\guillemotright{}\\\textbf{systemDiscovery}
        \nodepart{second} \texttt{register}\\ \texttt{lookup}\\ \texttt{revoke}};
  \end{tikzpicture}
  \caption{SysML-style block of the \texttt{systemDiscovery} service and its operations.}
  \label{fig:SD}
\end{figure}

The following interface operations are available.

\osubsection{register}{\mref{SystemRegisterRequest}}{\mref{SystemResponse}}

Registers the requesting system (name from its identity) with addresses, version, metadata and an optional device. Returns 201 for a new entry, 200 if an identical entry exists; an existing system with different data is refused (400).

\osubsection{lookup}{\mref{SystemLookupRequest}}{\mref{SystemListResponse}}

Returns the systems matching all given filters. With \texttt{verbose=\allowbreak{}true} (and \texttt{discovery.\allowbreak{}verbose} enabled) the device details are included.

Optional parameter: \texttt{verbose}.

\osubsection{revoke}{}{}

Removes the requesting system and its service instances. Returns 200 if removed, 204 if it was not registered.

\newpage

\section{Information Model}
\label{sec:model}

Here, all data objects that can be part of the \texttt{systemDiscovery} service are listed in alphabetic order.
Note that each subsection, which describes one type of object, begins with the \textit{struct} keyword, which is used to denote a collection of named fields, each with its own data type.
As a complement to the explicitly defined types in this section, there is also a list of implicit primitive types in Section \ref{sec:model:primitives}, which are used to represent things like hashes and identifiers.

\msubsection{struct}{Address}

A network address with its type.


\noindent\begin{tabularx}{\textwidth}{| >{\raggedright\arraybackslash}p{3.9cm} | >{\raggedright\arraybackslash\small}p{5.3cm} | X |} \hline
\rowcolor{gray!33} Field & Type & Description \\ \hline
type & \pref{String} & Address type: \texttt{HOSTNAME}, \texttt{IPV4}, \texttt{IPV6} or \texttt{MAC}. \\ \hline
address & \pref{String} & One network address (host name, IPv4, IPv6 or MAC address). \\ \hline
\end{tabularx}



\msubsection{struct}{DeviceResponse}

A registered device.


\noindent\begin{tabularx}{\textwidth}{| >{\raggedright\arraybackslash}p{3.9cm} | >{\raggedright\arraybackslash\small}p{5.3cm} | X |} \hline
\rowcolor{gray!33} Field & Type & Description \\ \hline
name & \pref{Name} & Device name. \\ \hline
metadata & \pref{Metadata} & Free-form metadata (JSON object). \\ \hline
addresses & \pref{List}$<$\allowbreak{}\mref{Address}$>$ & Network addresses with their types. \\ \hline
created\allowbreak{}At & \pref{DateTime} & UTC time of creation. \\ \hline
updated\allowbreak{}At & \pref{DateTime} & UTC time of the last change. \\ \hline
\end{tabularx}



\msubsection{struct}{MetadataRequirement}

A JSON object of metadata requirements: each key is a metadata path (dot notation), each value the required value or an operator object such as \texttt{\{"op": "LESS\_\allowbreak{}THAN",\allowbreak{} "value": 10\}}.



\msubsection{struct}{SystemListResponse}

List of systems.


\noindent\begin{tabularx}{\textwidth}{| >{\raggedright\arraybackslash}p{3.9cm} | >{\raggedright\arraybackslash\small}p{5.3cm} | X |} \hline
\rowcolor{gray!33} Field & Type & Description \\ \hline
entries & \pref{List}$<$\allowbreak{}\mref{SystemResponse}$>$ & Matching records of the current page. \\ \hline
count & \pref{Number} & Total number of matching records (all pages). \\ \hline
\end{tabularx}



\msubsection{struct}{SystemLookupRequest}

Filter for systems; all fields optional, combined with AND.


\noindent\begin{tabularx}{\textwidth}{| >{\raggedright\arraybackslash}p{3.9cm} | >{\raggedright\arraybackslash\small}p{5.3cm} | X |} \hline
\rowcolor{gray!33} Field & Type & Description \\ \hline
system\allowbreak{}Names & \pref{List}$<$\allowbreak{}\pref{Name}$>$ & System names to match. \\ \hline
addresses & \pref{List}$<$\allowbreak{}\pref{String}$>$ & Network addresses of the entity. \\ \hline
address\allowbreak{}Type & \pref{String} & Address type filter: \texttt{HOSTNAME}, \texttt{IPV4}, \texttt{IPV6} or \texttt{MAC}. \\ \hline
metadata\allowbreak{}Requirement\allowbreak{}List & \pref{List}$<$\allowbreak{}\mref{MetadataRequirement}$>$ & Metadata requirements; each map is one alternative. \\ \hline
versions & \pref{List}$<$\allowbreak{}\pref{Version}$>$ & System versions to match. \\ \hline
device\allowbreak{}Names & \pref{List}$<$\allowbreak{}\pref{Name}$>$ & Device names to match. \\ \hline
\end{tabularx}



\msubsection{struct}{SystemRegisterRequest}

Registers the requesting system; the name is taken from its authenticated identity.


\noindent\begin{tabularx}{\textwidth}{| >{\raggedright\arraybackslash}p{3.9cm} | >{\raggedright\arraybackslash\small}p{5.3cm} | X |} \hline
\rowcolor{gray!33} Field & Type & Description \\ \hline
metadata & \pref{Metadata} & Free-form metadata (JSON object). \\ \hline
version & \pref{Version} & Version of the system. \\ \hline
addresses & \pref{List}$<$\allowbreak{}\pref{String}$>$ & Network addresses of the entity. \\ \hline
device\allowbreak{}Name & \pref{Name} & Name of the device hosting the system. \\ \hline
\end{tabularx}



\msubsection{struct}{SystemResponse}

A registered system.


\noindent\begin{tabularx}{\textwidth}{| >{\raggedright\arraybackslash}p{3.9cm} | >{\raggedright\arraybackslash\small}p{5.3cm} | X |} \hline
\rowcolor{gray!33} Field & Type & Description \\ \hline
name & \pref{Name} & System name. \\ \hline
metadata & \pref{Metadata} & Free-form metadata (JSON object). \\ \hline
version & \pref{Version} & Version of the system. \\ \hline
addresses & \pref{List}$<$\allowbreak{}\mref{Address}$>$ & Network addresses with their types. \\ \hline
device & \mref{DeviceResponse} & The device hosting the system, if registered. \\ \hline
created\allowbreak{}At & \pref{DateTime} & UTC time of creation. \\ \hline
updated\allowbreak{}At & \pref{DateTime} & UTC time of the last change. \\ \hline
\end{tabularx}



\subsection{Primitives}
\label{sec:model:primitives}

Types and structures mentioned throughout this document that are assumed to be available to implementations of this service.
The concrete interpretations of each of these types and structures must be provided by any IDD document claiming to implement this service.

\noindent\begin{tabularx}{\textwidth}{| p{4.6cm} | X |} \hline
\rowcolor{gray!33} Type & Description \\ \hline
\pdef{DateTime} & Pinpoints a specific moment in time (UTC). \\ \hline
\pdef{List} & An \textit{array} of a known number of items, each having type A. \\ \hline
\pdef{Metadata} & A set of key-value pairs with arbitrary, possibly nested values. \\ \hline
\pdef{Name} & A string identifier that is intended to be both human and machine-readable. \\ \hline
\pdef{Number} & An integer number. \\ \hline
\pdef{String} & A chain of characters. \\ \hline
\pdef{Version} & Specifies a version, semantic versioning \texttt{major.minor.patch}. \\ \hline
\end{tabularx}

\newpage

\bibliographystyle{IEEEtran}
\bibliography{bibliography}

\newpage

\section{Revision History}
\subsection{Amendments}

\noindent\begin{tabularx}{\textwidth}{| p{1cm} | p{3cm} | p{2cm} | X | p{4cm} |} \hline
\rowcolor{gray!33} No. & Date & Version & Subject of Amendments & Author \\ \hline

1 & 2026-10-22 & \arrowversion & Document created from the Arrowhead LaTeX templates for the 5.2.2 Java Spring implementation, including the local extensions of this workspace. & Jerker Delsing \\ \hline
\end{tabularx}

\subsection{Quality Assurance}

\noindent\begin{tabularx}{\textwidth}{| p{1cm} | p{3cm} | p{2cm} | X |} \hline
\rowcolor{gray!33} No. & Date & Version & Approved by \\ \hline

1 & & \arrowversion  &  \\ \hline

\end{tabularx}

\end{document}
